\scanpage{188}
\mspage{II.36}
% archive sheet number: 187
% block: top, line 1 | display
\[
\pi^{\mathrm{alg}}=\deleted{\unsure{\operatorname{Lie}_\ell|\varprojlim}}\text{``}\varprojlim\text{''}\widetilde G_E
\]

% block: upper third, lines 2-3 | text
où $\widetilde G_E$ est l'\uline{enveloppe algébrique} du groupe linéaire $G_E$.

% block: upper third, lines 4-6 | text
\uline{NB} En fait, la flèche $\unsure{\alpha}$ est un \uline{isom}, car tt groupe analytique compact sur $\mathbb Q_\ell$ a une représentation linéaire fidèle\ldots
\margin{\unsure{réf?}}

% block: middle third, lines 7-10 | text
Donc $\pi^{\mathrm{an}}\longrightarrow\pi^{\mathrm{alg}}$ est un isom de \deleted{\unsure{groupes}}\inserted{pro-groupes} analytiques ssi \inserted{tous} les $G_E$ sont \inserted{tous} des groupes \deleted{\unsure{linéaires}} algébriques. \inserted{Ce n'est sans doute jamais le cas, sauf si $\pi\simeq\pi^{\mathrm{an}}$ (donc $\pi\longleftarrow\pi^{\mathrm{alg}}$)}
\margin{?}

% block: middle third, lines 11-13 | deletion
\deleted{\unsure{Il suffit p.ex. que ce soient} leurs propres groupes dérivés --- en particulier qu'ils soient semi simples\ldots}

% block: middle third, line 14 | text
L'hom (6.6.1) induit un hom

% block: middle to lower third, line 15 | display
\[
\mathfrak g=\operatorname{Lie}^{\mathrm{an}}(\pi)\unsure{\longrightarrow|\text{---}}\widetilde{\mathfrak g}=\operatorname{Lie}^{\mathrm{alg}}(\pi).
\tag*{(6.6.2)}
\]

% block: lower third, lines 16-19 | text
C'est un isomorphisme si et seulement si les Algèbres de Lie $\mathfrak g_E=\operatorname{Lie}(G_E)$ sont algébriques, -- p.ex si ce sont leurs propres Algèbres dérivées, i.e. les $G_E$ sont \deleted{d'indice}\inserted{ouverts} dans les $\widetilde G_E$. Dans ce cas, \deleted{\unsure{l'hom}} foncteur

% block: lower third, lines 20-21 | display
\[
\varinjlim_{\substack{U_i\text{ ouvert}\\\text{de }\pi}}\operatorname{Mod}_{\mathbb Q_\ell}(U_i)\simeq\operatorname{Mod}(\widetilde{\mathfrak g})
\tag*{(6.6.3)}
\]

% block: lower third, line 22 | text
est évidemment une équivalence de catégories\ldots.

% block: bottom left margin | margin note
% slanted
\margin{laisser tomber remarques 6.7. et 6.8. dans \unsure{un texte|au net|au verso}}

% block: bottom, lines 23-24 | text
\uline{Remarque} 6.7. Si on interprète le premier \unsure{mb|mbr} en (6.6.3) comme $\operatorname{Mod}(\pi^{\mathrm{alg}})$, l'équivalence
